\begin{diagram}[
  type=flow,
  label={fig:sec01-lifecycle},
  caption={The data lifecycle proceeds from sources to serving, but it is a feedback loop rather than a rigid one-way chain: quality and reliability, metadata and governance, observability, security, and cost are cross-cutting concerns that apply across every stage, not steps that happen only after orchestration.},
  source={Author's synthesis.},
  description={A top-down flow of six sequential stages in reading order: Sources, Ingestion, Storage, Processing and Transformation, Orchestration, and Serving. A dashed band drawn behind all six stages represents cross-cutting concerns -- quality and reliability, metadata and governance, observability, security, and cost -- that apply across every stage rather than following it as a seventh or eighth step.},
  minspace={0.52\textheight}
]
% Six stages are too dense for the default horizontal spacing, so this
% diagram places them directly with \RKNode/\flowedge in a vertical column,
% as before. Quality/reliability and metadata/governance were previously
% drawn as sequential stages 6 and 7, which wrongly implied they happen only
% after orchestration. They are cross-cutting (so are observability,
% security, and cost, which the caption already excluded from the chain), so
% they are drawn once, behind the whole chain, as a shaded dashed band with
% its own label rather than as ordinary flow nodes. The band is drawn first
% so the sequential nodes and edges render on top of it and its label sits
% clear of every node's text (see the x-offset below).
\node[
  fill=Muted!8,
  draw=Muted,
  dashed,
  thick,
  rounded corners=4pt,
  anchor=north west,
  minimum width=6.8cm,
  minimum height=11.05cm,
  inner sep=0pt
] at (-2.4,0.9) {};
\node[
  anchor=north,
  text width=2.4cm,
  align=left,
  font=\sffamily\fontsize{6.8}{8.1}\selectfont,
  text=Muted
] at (3.1,-4.2) {\textbf{Cross-cutting, every stage:}\\[3pt]
Quality and\\Reliability\\[3pt]
Metadata and\\Governance\\[3pt]
Observability\\[3pt]
Security\\[3pt]
Cost};
\RKNode{sources}{0}{0}{Sources}
\RKNode{ingestion}{0}{-1.85}{Ingestion}
\RKNode{storage}{0}{-3.70}{Storage}
\RKNode{processing}{0}{-5.55}{Processing and Transformation}
\RKNode{orchestration}{0}{-7.40}{Orchestration}
\RKNode{serving}{0}{-9.25}{Serving}
\flowedge{sources}{ingestion}
\flowedge{ingestion}{storage}
\flowedge{storage}{processing}
\flowedge{processing}{orchestration}
\flowedge{orchestration}{serving}
\end{diagram}
